% cuadro C2_1
\begin{table}[htbp]
\centering
\scriptsize
\caption{Ingresos por renglón del artículo 1o.: iniciativa 2025 contra iniciativa 2026 (línea P)}
\label{cua:C2_1}
\begin{tabular}{p{5.4cm}rrrrrr}
\toprule
\textbf{Concepto} & \textbf{ILIF 2025} & \textbf{ILIF 2026} & \textbf{Dif. nominal} & \textbf{Dif. real} & \textbf{Var. real \%} & \textbf{\% PIB 2026} \\
\midrule
1. Impuestos & 5 297 812,9 & 5 838 541,1 & 540 728,2 & 286 433,2 & 5,16 & 15,08 \\
   11.01 Impuesto sobre la renta & 2 859 575,1 & 3 070 149,1 & 210 574,0 & 73 314,4 & 2,45 & 7,93 \\
   13.01 Impuesto al valor agregado & 1 463 279,9 & 1 589 069,0 & 125 789,1 & 55 551,7 & 3,62 & 4,10 \\
   13.02 Impuesto especial sobre produccion y servicios & 713 844,0 & 761 501,9 & 47 657,9 & 13 393,4 & 1,79 & 1,97 \\
   14 Impuestos al comercio exterior & 151 789,7 & 254 756,8 & 102 967,1 & 95 681,2 & 60,15 & 0,66 \\
   17 Accesorios de impuestos & 81 497,4 & 135 769,4 & 54 272,0 & 50 360,1 & 58,96 & 0,35 \\
2. Cuotas y aportaciones de seguridad social & 603 077,9 & 641 782,1 & 38 704,2 & 9 756,5 & 1,54 & 1,66 \\
4. Derechos & 137 500,5 & 157 081,7 & 19 581,2 & 12 981,2 & 9,01 & 0,41 \\
5. Productos & --- & 16 488,3 & --- & --- & --- & 0,04 \\
6. Aprovechamientos & 223 166,3 & 203 520,5 & -19 645,8 & -30 357,8 & -12,98 & 0,53 \\
7. Ingresos por ventas de bienes y prestacion de servicios & 1 500 579,0 & 1 630 973,6 & 130 394,6 & 58 366,8 & 3,71 & 4,21 \\
9.97 Transferencias del Fondo Mexicano del Petroleo & 279 766,8 & 232 630,4 & -47 136,4 & -60 565,2 & -20,66 & 0,60 \\
0. Ingresos derivados de financiamientos & 1 246 366,5 & 1 472 626,4 & 226 259,9 & 166 434,3 & 12,74 & 3,80 \\
\midrule
\textbf{TOTAL} & \textbf{9 302 015,8} & \textbf{10 193 683,7} & \textbf{891 667,9} & \textbf{445 171,1} & \textbf{4,57} & \textbf{26,33} \\
\bottomrule
\end{tabular}
\\[2pt]\begin{minipage}{0.92\textwidth}\scriptsize Línea P: ex ante contra ex ante. Ambas columnas son iniciativa, no ley aprobada. Millones de pesos corrientes; la diferencia real usa el deflactor del PIB de 4.8 por ciento que declara el CGPE 2026. Fuente: ILIF 2026 e ILIF 2025, artículo 1o. Elaborado por el ITED.\end{minipage}
\end{table}
